\documentclass[10pt]{article}

\usepackage{latexsym}
\usepackage{amssymb}
\usepackage{code}
\usepackage{graphicx}

\usepackage{listings}



\usepackage[english]{babel}
%\usepackage{blindtext}
%\usepackage{fontspec}
%\setmainfont{Arial}

%\renewcommand{\familydefault}{\sfdefault}
%\usepackage{blindtext}
\pagenumbering{arabic}  % Arabic page numbers GM July 2000

\usepackage{authoraftertitle}
\usepackage{fancyhdr}
% Clear the header and footer
\fancyhead{}
\fancyfoot{}
%\lhead{\includegraphics[height=0.7cm]{../logo/nyp-logo-int.png} }
\rhead{\scriptsize \MyTitle}
%\lfoot{\scriptsize Specialist Diploma in Business \& Big Data Analytics
%  \\ Copyright
 % \copyright\ Nanyang Polytechnic. All Rights Reserved.}
\rfoot{\thepage}
\pagestyle{fancy}


\renewcommand{\topfraction}{0.95}
\renewcommand{\textfraction}{0.02}
\renewcommand{\floatpagefraction}{0.95}


\bibliographystyle{plain}

%%\textwidth      150mm \textheight     210mm \oddsidemargin  -2mm \evensidemargin -2mm

\renewcommand{\baselinestretch}{0.987}

\setlength{\parskip}{0.0in}

%\input{macros-ms}

\newcommand{\Nturns}{\, \vdash_{\mbox{\scriptsize lnf}} \,}

\newcommand{\tr}[1]{}
%%\newcommand{\tr}[1]{#1}
%%\newcommand{\nottr}[1]{#1}
\newcommand{\nottr}[1]{}

%\newcommand{\implies}{\supset}
\newcommand{\clabel}[1]{\mbox{(#1)}}
\newcommand{\rat}[1]{\rightarrowtail_{#1}}
\newcommand{\arr}{\rightarrow}
\newcommand{\arrow}{\rightarrow}
\newcommand{\Arr}{\Rightarrow}
\newcommand{\atsign}{@}
\newcommand{\simparrow}[0]{\Longleftrightarrow}
\newcommand{\proparrow}[0]{\Longrightarrow}
\newcommand{\comment}[1]{}
\newcommand{\ignore}[1]{}
%\newcommand{\kl}[1]{{\bf KL:#1}} %%%{\marginpar{\sc kl}{\bf #1}}
\newcommand{\kl}[1]{}
\newcommand{\ms}[1]{{\bf MS:#1}}       %%{\marginpar{\sc ms}{\bf #1}}
%%\newcommand{\jw}[1]{\marginpar{\sc jw}{\bf #1}}
\newcommand{\pjs}[1]{}
%%\newcommand{\ms}[1]{}
\newcommand{\jw}[1]{}
%%\newcommand{\kl}[1]{}


\newcommand{\pow}{\^{}}
\newcommand{\venv}{\Delta}
\newcommand{\mleq}{\mbox{\tt leq}}
\newcommand{\mmleq}{\leq}
\newcommand{\mas}{\mbox{\tt as}}
\newcommand{\mfix}{\mu}

\newcommand{\mysection}[1]{\vspace*{-2mm}\section{#1}\vspace*{-1mm}}
\newcommand{\mysubsection}[1]{\vspace*{-1mm}\subsection{#1}\vspace*{-1mm}}

\newcounter{cnt}
\newtheorem{ex}{Example}
%\newenvironment{example}{
%        \begin{ex}\rm}%
%    {\hfill$\Box$\end{ex}}

\newenvironment{nexample}{
        \begin{ex}\rm}%
        {\end{ex}}

%%\newtheorem{rem}{Remark}
%%\renewenvironment{remark}{
%%        \begin{rem}\rm}%
%%    {\hfill$\Box$\end{rem}}

%%\newenvironment{myexample}{
%%        \begin{ex}\rm}%
%%  {\hfill $\Diamond$\end{ex}}
%\newenvironment{example}{
%        \begin{example}\rm}%
%    {\end{example}}

\newenvironment{ttline}{\begin{trivlist}\item \tt}{\end{trivlist}}
\newenvironment{ttprog}{\begin{trivlist}\small\item \tt
        \begin{tabbing}}{\end{tabbing}\end{trivlist}}


\newcommand{\figcode}[1]
        {\begin{figure*}[t]#1
        \end{figure*}}


\title{Introduction to Spark Machine Learning}
%\author{Kenny Zhuo Ming Lu\\
%  \multicolumn{1}{p{.7\textwidth}}{\centering
%  \emph{School of Information Technology \\ Nanyang Polytechnic \\
%    180 Ang Mo Kio Avenue 8, Singapore 569830}}
%}


\begin{document}
\maketitle \makeatactive
\thispagestyle{fancy}

%\lstset{language=scala}

%\begin{abstract}
%\end{abstract}



\section{Exercises }

\begin{enumerate}
   \item Start the Spark cluster
\begin{code}
$/opt/spark-1.5.2-bin-hadoop2-hive2-r/sbin/start-all.sh
\end{code}
%$
   \item We are going to transform some data from HDFS. 
   \item Start Hadoop. In the VM, start a terminal and run the following command
\begin{code}
$ start-dfs.sh && start-yarn.sh
\end{code} %$
   \item Check out the source codes.
           \begin{enumerate}
               \item If you find the {\tt \~{}/git/scala-bigdata-starter/} folder exists in your system, 
\begin{code}
$ cd ~/git/scala-bigdata-starter
$ git pull
\end{code}
               \item otherwise, clone it from github.
\begin{code}
$ cd ~
$ mkdir git
$ cd git 
$ git clone http://github.com/luzhuomi/scala-bigdata-starter.git
$ cd scala-bigdata-starter
\end{code}
%$
             \end{enumerate}
 \item Change to the following folder
\begin{code}
$ cd ~/git/scala-bigdata-starter/spark-examples/
\end{code}
%$

\end{enumerate}

\subsection{Tweet classifier}

The school have collected some twitter data on Joint Admission
Exercise (in Singapore). They are often mentioned in the short form
``jae''. However the twitter data collector can only specify the
keyword ``jae'', we are looking the admision exercise related data yet
we are observing data talking about KPOP. (Because there are many
Korean celebrities having ``jae'' as part of their names.) The task
here is to build a classifier to differentiate the KPOP tweets
mentioning ``jae'' as a person's name or otherwise.




\end{document}
